\documentclass[11pt]{article}
\usepackage[margin=1in]{geometry}
\usepackage[T1]{fontenc}
\usepackage[utf8]{inputenc}
\usepackage{lmodern}
\usepackage{microtype}
\usepackage{array}
\usepackage{booktabs}
\usepackage{enumitem}
\usepackage{hyperref}
\usepackage{xcolor}
\hypersetup{
  colorlinks=true,
  linkcolor=black,
  urlcolor=blue!50!black,
  pdftitle={Coding Scheme Dossier}
}
\setlist[itemize]{leftmargin=1.6em}
\setlist[description]{style=nextline,leftmargin=3.5cm,labelsep=0.6em}
\newcommand{\field}[1]{\nolinkurl{#1}}
\newcommand{\filepath}[1]{\nolinkurl{#1}}
\newcommand{\schemeheader}[3]{
  \begin{center}
    {\LARGE \textbf{#1}}\\[0.5em]
    {\large #2}\\[0.4em]
    {\normalsize \textit{#3}}
  \end{center}
  \vspace{0.8em}
}



\begin{document}
\schemeheader
  {Scheme 3}
  {Stage 3 Full-Text Deep Coding Scheme}
  {Musculoskeletal-focused full-text adjudication and interpretive coding}

\section*{Purpose}
This scheme is the full-text deep coding framework used for Stage 3 candidate reviews, especially musculoskeletal chronic pain reviews. It captures conceptual depth that cannot be resolved reliably at the abstract level: coverage of each BPS domain, pairwise and triadic integration quality, biopsychosocial typology, psychological concepts, theoretical frameworks, conceptual problems, and evidential quotations.

\section*{Operational Status}
\begin{itemize}
\item Workflow position: Stage 3 full-text coding after candidate identification and retrieval triage.
\item Operational mode: human-coded template with pilot and reliability subsamples; protocol allows AI assistance for concept mapping, but final adjudication remains human.
\item Canonical implementation basis: current full-text template and prose codebook.
\end{itemize}

\section*{Canonical Source Paths}
\begin{itemize}
\item \filepath{src/protocol/codebooks/stage3_codebook.md}
\item \filepath{src/review_stages/04_extraction/codebooks/stage3_codebook.csv}
\item \filepath{src/bps_review/extraction/stage3_prep.py}
\item \filepath{src/review_stages/04_extraction/forms/stage3_fulltext_coding_template.csv}
\item \filepath{src/review_stages/04_extraction/forms/stage3_pilot_sample.csv}
\item \filepath{src/review_stages/04_extraction/forms/stage3_reliability_sample.csv}
\end{itemize}

\section*{Core Fields and Controlled Values}
\begin{description}
\item[\field{record_id}] Internal record identifier.
\item[\field{full_text_available}] \texttt{yes}; \texttt{no}; \texttt{partial}.
\item[\field{pain_condition_detail}] Free-text specification of the pain condition studied in the review.
\item[\field{domain_coverage_bio}] \texttt{elaborated}; \texttt{mentioned}; \texttt{minimal}; \texttt{absent}.
\item[\field{domain_coverage_psych}] \texttt{elaborated}; \texttt{mentioned}; \texttt{minimal}; \texttt{absent}.
\item[\field{domain_coverage_social}] \texttt{elaborated}; \texttt{mentioned}; \texttt{minimal}; \texttt{absent}.
\item[\field{integration_bio_psych}] \texttt{mechanistic}; \texttt{directional}; \texttt{descriptive}; \texttt{mentioned}; \texttt{none}.
\item[\field{integration_psych_social}] \texttt{mechanistic}; \texttt{directional}; \texttt{descriptive}; \texttt{mentioned}; \texttt{none}.
\item[\field{integration_bio_social}] \texttt{mechanistic}; \texttt{directional}; \texttt{descriptive}; \texttt{mentioned}; \texttt{none}.
\item[\field{integration_triadic}] \texttt{mechanistic}; \texttt{descriptive}; \texttt{partial}; \texttt{none}.
\item[\field{integration_mechanism_summary}] Concise free-text summary of proposed cross-domain pathways.
\item[\field{overall_balance}] \texttt{balanced}; \texttt{psych-dominant}; \texttt{bio-dominant}; \texttt{social-dominant}; \texttt{dyadic}; \texttt{unclear}.
\item[\field{bps_typology}] \field{true_integrative}; \field{multifactorial}; \field{pseudo_bps}; \field{rhetorical_bps}; \field{narrow_despite_label}; \field{unclear}.
\item[\field{concept_definitions_present}] \texttt{yes}; \texttt{partial}; \texttt{no}.
\item[\field{psychological_concepts_fulltext}] Full-text concept list, normally normalized and semicolon-delimited.
\item[\field{theoretical_frameworks_fulltext}] Full-text framework list, normally normalized and semicolon-delimited.
\item[\field{conceptual_problems_fulltext}] Conceptual issues such as vague definitions, construct overlap, tokenistic BPS use, missing social analysis, missing biology, mechanistic absence, unclear boundaries, or other.
\item[\field{integration_quotes_or_evidence}] Supporting quotations, section references, or evidential anchors from the full text.
\item[\field{coder_id} / \field{coder_notes} / \field{adjudication_status}] Provenance and adjudication tracking fields.
\end{description}

\section*{Typology Interpretation Rules}
\begin{description}
\item[\field{true_integrative}] Explicit cross-domain causal or mechanistic interaction is central to the review's logic.
\item[\field{multifactorial}] Multiple domains are covered meaningfully but mostly in parallel rather than as an integrated explanatory model.
\item[\field{pseudo_bps}] BPS language is used, but one or more core domains are thin, absent, or tokenistic.
\item[\field{rhetorical_bps}] BPS is used mainly as framing, justification, or rhetorical positioning without substantive analytic integration.
\item[\field{narrow_despite_label}] The review claims a BPS framing, but substantive scope is essentially single-domain or severely restricted.
\end{description}

\section*{Coding Expectations}
\begin{itemize}
\item Domain coverage is coded at the level of substantive treatment, not keyword presence alone.
\item Pairwise integration codes should reflect whether the review articulates mechanisms, directional influences, descriptive links, or only passing mention.
\item Triadic integration should be reserved for genuinely three-domain reasoning rather than serial mention of separate domains.
\item The scheme is designed to combine quantitative tallying with qualitative excerpting so that research questions on both distribution and interpretation can be answered from the same coding file.
\end{itemize}

\section*{Reliability Architecture}
\begin{itemize}
\item Pilot sample: up to 5 records generated through \filepath{src/bps_review/extraction/stage3_prep.py}.
\item Reliability sample: roughly 20\% of candidates, capped operationally in the form-generation logic.
\item Reliability fields include independent reviewer ratings for domain coverage, triadic integration, BPS typology, and free-text notes, followed by adjudicated values.
\end{itemize}

\section*{Implementation Notes and Documented Divergences}
\begin{itemize}
\item The prose Stage 3 codebook includes \texttt{partial} and \texttt{mentioned} intermediary integration values more explicitly than the CSV codebook; this dossier follows the prose codebook plus the actual form field names because that is the richer operational specification.
\item The CSV codebook uses shorter names such as \field{bio_coverage}; the current generated template uses \field{domain_coverage_bio}. The template naming is treated here as canonical because it is the actual working form.
\end{itemize}

\section*{Primary Outputs and Forms Using This Scheme}
\begin{itemize}
\item \filepath{src/review_stages/04_extraction/forms/stage3_fulltext_coding_template.csv}
\item \filepath{src/review_stages/04_extraction/forms/stage3_pilot_sample.csv}
\item \filepath{src/review_stages/04_extraction/forms/stage3_reliability_sample.csv}
\item \filepath{src/review_stages/04_extraction/outputs/stage3_candidate_manifest.csv}
\end{itemize}

\end{document}
